\documentclass[12pt,a4paper]{article}

% --- Packages ---
\usepackage[utf8]{inputenc}
\usepackage[margin=1in]{geometry}
\usepackage{amsmath,amssymb,amsfonts,amsthm}
\usepackage{graphicx}
\usepackage{xcolor}
\usepackage{hyperref}
\usepackage{listings}
\usepackage{booktabs}
\usepackage{tikz}
\usetikzlibrary{shapes.geometric, arrows, positioning, shadows, fit, backgrounds, patterns, calc}
\usepackage{float}
\usepackage{caption}
\usepackage{subcaption}
\usepackage{enumitem}

% --- Hyperref Setup ---
\hypersetup{
    colorlinks=true,
    linkcolor=blue!80!black,
    citecolor=blue!80!black,
    urlcolor=blue!80!black,
    pdftitle={Methodology: Camera Control and Data Acquisition Pipeline},
    pdfauthor={System Documentation}
}

% --- Custom Code Listing Styling ---
\lstset{
    backgroundcolor=\color{gray!5!white},
    basicstyle=\ttfamily\small,
    breaklines=true,
    captionpos=b,
    commentstyle=\color{green!50!black},
    keywordstyle=\color{blue!80!black},
    stringstyle=\color{red!70!black},
    frame=single,
    numbers=left,
    numberstyle=\tiny\color{gray},
    stepnumber=1,
    tabsize=4
}

\begin{document}

\section{Methodology}

This section describes the software architecture, USB Picture Transfer Protocol (PTP/ISO 15740) data acquisition pipeline, and system modeling for the Canon EOS R8 FastAPI control system.

\subsection{System Data Flow Architecture}

The system connects a web-based user interface to physical camera hardware via an asynchronous FastAPI backend and low-level \texttt{libgphoto2} PTP drivers. The architecture consists of four primary operational layers:

\begin{enumerate}[label=\arabic*.]
    \item \textbf{Frontend UI (Browser):} HTML5/JavaScript dashboard managing live preview displays and user interaction controls.
    \item \textbf{Backend API (FastAPI):} Asynchronous web application handling routing, CORS, HTTP response streams, and endpoint contracts.
    \item \textbf{Controller Layer (\texttt{CameraController}):} Python thread-synchronized wrapper interfacing with \texttt{python-gphoto2} bindings and managing hardware locks.
    \item \textbf{Hardware \& Protocol Layer (\texttt{libgphoto2} / USB PTP):} USB bulk endpoints transferring operational data packets to/from Canon EOS R8 SDRAM and CMOS sensor assemblies.
\end{enumerate}

\begin{figure}[H]
\centering
\tikzset{
    block/.style = {rectangle, draw=blue!70!black, fill=blue!8, text width=7.2em, text centered, rounded corners=4pt, minimum height=2.6em, font=\bfseries\small},
    ctrlblock/.style = {rectangle, draw=orange!70!black, fill=orange!15, text width=9.5em, text centered, rounded corners=4pt, minimum height=2.6em, font=\bfseries\small},
    line/.style = {draw, -latex, thick, color=black!80}
}
\begin{tikzpicture}
    % Row 1: UI (x=0.8, y=2.4) <---> FastAPI (x=6.4, y=2.4)
    \node [block, fill=green!15, draw=green!60!black] (ui)  at (0.8, 2.4) {Browser Frontend \\ (\texttt{main.js})};
    \node [block, fill=blue!15, draw=blue!60!black]   (api) at (6.4, 2.4) {FastAPI Server \\ (\texttt{app.py})};
    
    % Row 2: Camera (x=0.8, y=0.0) <---> Controller (x=6.4, y=0.0)
    \node [block, fill=red!15, draw=red!60!black]     (cam)  at (0.8, 0.0) {Canon EOS R8 \\ Sensor / RAM};
    \node [ctrlblock]                                (ctrl) at (6.4, 0.0) {CameraController \\ (\texttt{camera\_ctl.py})};

    % Top Connections (UI <-> API)
    \draw [line] ($(ui.east)+(0,0.18)$) -- node[above, pos=0.5, font=\tiny\bfseries, fill=white, inner sep=1.5pt] {HTTP REST} ($(api.west)+(0,0.18)$);
    \draw [line] ($(api.west)-(0,0.18)$) -- node[below, pos=0.5, font=\tiny\bfseries, fill=white, inner sep=1.5pt] {JSON Stream} ($(ui.east)-(0,0.18)$);

    % Vertical Connections (API <-> Controller)
    \draw [line] ($(api.south)-(0.25,0)$) -- node[left, pos=0.5, font=\tiny\bfseries, fill=white, inner sep=1.5pt] {Method Call} ($(ctrl.north)-(0.25,0)$);
    \draw [line] ($(ctrl.north)+(0.25,0)$) -- node[right, pos=0.5, font=\tiny\bfseries, fill=white, inner sep=1.5pt] {RAM Buffer} ($(api.south)+(0.25,0)$);

    % Bottom Connections (Camera <-> Controller)
    \draw [line] ($(ctrl.west)+(0,0.18)$) -- node[above, pos=0.5, font=\tiny\bfseries, fill=white, inner sep=1.5pt] {USB PTP} ($(cam.east)+(0,0.18)$);
    \draw [line] ($(cam.east)-(0,0.18)$) -- node[below, pos=0.5, font=\tiny\bfseries, fill=white, inner sep=1.5pt] {Raw Payload} ($(ctrl.west)-(0,0.18)$);
\end{tikzpicture}
\caption{System architecture and data flow topology with enlarged controller block.}
\label{fig:arch_topology}
\end{figure}


\subsection{UML Use Case Diagram}

Figure~\ref{fig:uml_usecase} illustrates the UML Use Case Diagram modeled after camera control software patterns. Actors are placed strictly outside the system boundary, and all use cases are cleanly separated into functional columns.

\begin{figure}[H]
\centering
\tikzset{
    actor/.style = {rectangle, draw=gray!80, fill=yellow!15, rounded corners=4pt, minimum width=2.3cm, minimum height=1.0cm, align=center, font=\bfseries\small},
    usecase/.style = {ellipse, draw=blue!80!black, fill=blue!5, thick, text width=2.0cm, text centered, minimum width=2.1cm, minimum height=0.8cm, font=\scriptsize},
    system/.style = {draw=black!80, thick, rectangle, fill=gray!3},
    rel/.style = {draw, thick, color=black!80},
    inc/.style = {draw, dashed, ->, >=latex, thick, color=red!70!black}
}
\begin{tikzpicture}

    % 1. System Boundary Box (x from -1.4 to 7.4, y from -0.5 to 5.6)
    \draw [system] (-1.4, -0.5) rectangle (7.4, 5.6);
    \node [font=\bfseries\small] at (3.0, 5.1) {System Boundary: Camera Control System};

    % 2. Use Cases Grid
    % Left Column (x = -0.2)
    \node [usecase] (uc_connect) at (-0.2, 4.3) {Connect \\ Session};
    \node [usecase] (uc_config)  at (-0.2, 2.2) {Configure \\ Exposure};
    \node [usecase] (uc_preview) at (-0.2, 0.1) {Live View \\ Stream};

    % Right Column (x = 6.2)
    \node [usecase] (uc_capture) at (6.2, 4.3) {Capture \\ Photo};
    \node [usecase] (uc_fetch)   at (6.2, 0.1) {Fetch USB \\ Payload};

    % 3. Actors OUTSIDE System Box (Linux V4L2 moved lower to y = -1.4 for a clear arrow path)
    \node [actor] (user)   at (-4.2, 2.2)  {User / Web UI \\ (Primary Actor)};
    \node [actor] (camera) at (10.2, 3.2)  {Canon EOS R8 \\ (Hardware Actor)};
    \node [actor] (v4l2)   at (10.2, -1.4) {Linux V4L2 \\ (Kernel Actor)};

    % 4. Actor Associations
    \draw [rel] (user.east) -- (uc_connect.west);
    \draw [rel] (user.east) -- (uc_config.west);
    \draw [rel] (user.east) -- (uc_preview.west);

    \draw [rel] (uc_connect.east) -- (camera.west);
    \draw [rel] (uc_capture.east) -- (camera.west);
    \draw [rel] (uc_fetch.east)   -- (camera.west);

    % Clear diagonal path from Live View Stream to Linux V4L2 passing cleanly below Fetch USB Payload
    \draw [rel] (uc_preview.east) -- (v4l2.west);

    % 5. Include Relationships
    \draw [inc] (uc_capture)  -- node[above, pos=0.5, font=\tiny, fill=white, inner sep=1.5pt] {\texttt{<<include>>}} (uc_connect);
    \draw [inc] (uc_capture)  -- node[right, pos=0.5, font=\tiny, fill=white, inner sep=1.5pt] {\texttt{<<include>>}} (uc_fetch);
    \draw [inc] (uc_config)   -- node[left, pos=0.5, font=\tiny, fill=white, inner sep=1.5pt] {\texttt{<<include>>}} (uc_connect);
    \draw [inc] (uc_preview)  -- node[left, pos=0.5, font=\tiny, fill=white, inner sep=1.5pt] {\texttt{<<include>>}} (uc_config);

\end{tikzpicture}
\caption{UML Use Case Diagram with clear actor association paths.}
\label{fig:uml_usecase}
\end{figure}


\subsection{UML Sequence Diagram: Data Acquisition Lifecycle}

Figure~\ref{fig:uml_sequence} details the **UML Sequence Diagram** for full-resolution data acquisition across all four system lifelines. All lifeline connections begin strictly at the bottom boundary of the entity header boxes, and side notes are offset to prevent lines from passing inside text.

\begin{figure}[H]
\centering
\tikzset{
    lifeline/.style = {draw=black!50, dashed, thick},
    entity/.style = {rectangle, draw=blue!80!black, fill=blue!10, rounded corners=3pt, minimum width=2.0cm, minimum height=0.85cm, font=\bfseries\small, align=center},
    msg/.style = {draw, ->, >=latex, thick, color=blue!80!black},
    ret/.style = {draw, dashed, ->, >=latex, thick, color=green!60!black},
    self/.style = {draw, ->, >=latex, thick, color=orange!80!black},
    box/.style = {draw=gray!80, fill=yellow!10, rounded corners=3pt, font=\tiny\ttfamily, align=left}
}
\begin{tikzpicture}

    % Lifeline X Positions (Total Width = 11.4cm, Gap = 3.8cm)
    \coordinate (C_client) at (0.0, 0);
    \coordinate (C_app)    at (3.8, 0);
    \coordinate (C_ctrl)   at (7.6, 0);
    \coordinate (C_cam)    at (11.4, 0);

    % Lifeline Header Nodes (Named for clean boundary connection)
    \node [entity] (head_ui)   at (C_client) {Browser UI};
    \node [entity] (head_app)  at (C_app)    {FastAPI Server};
    \node [entity] (head_ctrl) at (C_ctrl)   {CameraController};
    \node [entity] (head_cam)  at (C_cam)    {Canon Camera};

    % Vertical Lifeline Lines - CONNECT ONLY TO BOTTOM BOUNDARY OF HEADER BOXES
    \draw [lifeline] (head_ui.south)   -- (0.0, -19.4);
    \draw [lifeline] (head_app.south)  -- (3.8, -19.4);
    \draw [lifeline] (head_ctrl.south) -- (7.6, -19.4);
    \draw [lifeline] (head_cam.south)  -- (11.4, -19.4);

    % Step 1: HTTP POST Capture
    \draw [msg] (0.0, -1.2) -- node[above, pos=0.45, font=\scriptsize, fill=white, inner sep=1.5pt] {1. \texttt{POST /api/capture}} (3.8, -1.2);

    % Step 2: FastAPI calls Controller
    \draw [msg] (3.8, -2.6) -- node[above, pos=0.45, font=\scriptsize, fill=white, inner sep=1.5pt] {2. \texttt{capture\_image()}} (7.6, -2.6);

    % Step 3: Acquire Mutex Lock (Self Loop)
    \draw [self] (7.6, -3.7) -- ++(0.6,0) -- ++(0,-0.8) -- node[right, pos=0.5, font=\tiny, fill=white, inner sep=1pt] {3. Acquire \texttt{self.lock}} (7.6, -4.5);

    % Step 4: Disable Viewfinder over PTP
    \draw [msg] (7.6, -5.7) -- node[above, pos=0.45, font=\scriptsize, fill=white, inner sep=1.5pt] {4. \texttt{set\_config(viewfinder=0)}} (11.4, -5.7);

    % Step 5: ACK Viewfinder Disabled
    \draw [ret] (11.4, -7.0) -- node[below, pos=0.45, font=\tiny, fill=white, inner sep=1.5pt] {ACK (Viewfinder Disabled)} (7.6, -7.0);

    % Step 6: Hardware Settling Delay (Self Loop)
    \draw [self] (7.6, -8.1) -- ++(0.6,0) -- ++(0,-0.8) -- node[right, pos=0.5, font=\tiny, fill=white, inner sep=1pt] {5. \texttt{time.sleep(0.5s)}} (7.6, -8.9);

    % Step 7: Trigger Shutter Capture
    \draw [msg] (7.6, -10.1) -- node[above, pos=0.45, font=\scriptsize, fill=white, inner sep=1.5pt] {6. \texttt{capture(GP\_CAPTURE\_IMAGE)}} (11.4, -10.1);
    
    % Side Note Box - OFFSET TO THE RIGHT so vertical lifeline connects to boundary, not inside text
    \node [box, anchor=west] at (11.65, -10.9) {Sensor Exposes $\rightarrow$\\SDRAM Buffer};
    \draw [draw=black!50, dotted] (11.4, -10.9) -- (11.65, -10.9);

    % Step 8: USB Bulk Buffer Transfer
    \draw [msg] (7.6, -12.1) -- node[above, pos=0.45, font=\scriptsize, fill=white, inner sep=1.5pt] {7. \texttt{file\_get(camera\_file)}} (11.4, -12.1);
    \draw [ret] (11.4, -13.5) -- node[below, pos=0.45, font=\scriptsize, fill=white, inner sep=1.5pt] {Raw Binary Payload over USB} (7.6, -13.5);

    % Step 9: Re-enable Viewfinder
    \draw [msg] (7.6, -15.0) -- node[above, pos=0.45, font=\tiny, fill=white, inner sep=1.5pt] {8. \texttt{set\_config(viewfinder=1)}} (11.4, -15.0);

    % Step 10: Return File Path to FastAPI
    \draw [ret] (7.6, -16.3) -- node[above, pos=0.45, font=\scriptsize, fill=white, inner sep=1.5pt] {Target File Path} (3.8, -16.3);

    % Step 11: Return HTTP Response to Client
    \draw [ret] (3.8, -17.6) -- node[above, pos=0.45, font=\scriptsize, fill=white, inner sep=1.5pt] {9. HTTP 200 JSON Response} (0.0, -17.6);

\end{tikzpicture}
\caption{UML Sequence Diagram with boundary-connected lifelines and offset note box.}
\label{fig:uml_sequence}
\end{figure}


\subsection{PTP Data Extraction and Preview Mechanics}

\subsubsection{PTP Capture Target and In-Memory Storage}
Upon connection initialization, the controller sets \texttt{capturetarget = 1}. This directs the DIGIC X processor to store image data directly in volatile camera SDRAM buffer space rather than flushing to physical SD card memory, preventing file system locks.

\subsubsection{C-Python Buffer Pointer Retrieval}
The binary transfer is executed cleanly without file system overhead:
\begin{lstlisting}[language=Python]
# Allocate raw C memory file container wrapper
camera_file = gp.CameraFile()

# Fetch binary buffer directly over USB Bulk IN Endpoint
self.camera.file_get(file_path.folder, file_path.name, 
                     gp.GP_FILE_TYPE_NORMAL, camera_file, self.context)

# Flush binary data buffer to target path
camera_file.save(target_path)
\end{lstlisting}

\subsubsection{Live Preview Frame Extraction}
For real-time streaming, the backend continuously pulls preview frames from the viewfinder buffer memory via \texttt{capture\_preview()}:
\begin{lstlisting}[language=Python]
def capture_preview(self):
    with self.lock:
        camera_file = gp.CameraFile()
        self.camera.capture_preview(camera_file, self.context)
        file_data = camera_file.get_data_and_size()
        return memoryview(file_data).tobytes()
\end{lstlisting}
The FastAPI generator packages these frame buffers into an HTTP stream (\texttt{multipart/x-mixed-replace; boundary=frame}) refreshed at $\approx 20\,\text{FPS}$.

\end{document}
